\subsection{Feedforward Neural Networks}
A Feedforward Neural Network (FFNN) is a type of artificial neural network that define a mapping from the input space to an output space. 
The learned mapping $f(\vct{x})$ is defined as a composition of several functions, each representing a layer of the network.
\begin{figure}[H]
    \centering
    \scalebox{0.8}{
        \begin{tikzpicture}[
  input/.style={circle, draw=teal!80!black, fill=teal!25, minimum size=22pt},
  hidden/.style={circle, draw=orange!80!black, fill=orange!25, minimum size=22pt},
  output/.style={circle, draw=violet!80!black, fill=violet!25, minimum size=22pt},
  layer label/.style={font=\small\bfseries, text=black!70, align=center},
  conn/.style={->, thick, black!30, opacity=0.6}
]

  % Input layer
  \foreach \i in {1,...,4} {
    \node[input] (i\i) at (0, -\i*1.1) {};
  }

  % Hidden layer 1
  \foreach \j in {1,...,3} {
    \node[hidden] (h1\j) at (1.8, -0.55 - \j*1.1) {};
  }

  % Hidden layer 2
  \foreach \j in {1,...,3} {
    \node[hidden] (h2\j) at (3.6, -0.55 - \j*1.1) {};
  }

  % Hidden layer 3
  \foreach \j in {1,...,3} {
    \node[hidden] (h3\j) at (5.4, -0.55 - \j*1.1) {};
  }

  % Output layer
  \foreach \k in {1,2} {
    \node[output] (o\k) at (7.2, -1.1 - \k*1.1) {};
  }

  % Connections: Input → Hidden 1
  \foreach \i in {1,...,4} {
    \foreach \j in {1,...,3} {
      \draw[conn] (i\i) -- (h1\j);
    }
  }

  % Connections: Hidden 1 → Hidden 2
  \foreach \i in {1,...,3} {
    \foreach \j in {1,...,3} {
      \draw[conn] (h1\i) -- (h2\j);
    }
  }

  % Connections: Hidden 2 → Hidden 3
  \foreach \i in {1,...,3} {
    \foreach \j in {1,...,3} {
      \draw[conn] (h2\i) -- (h3\j);
    }
  }

  % Connections: Hidden 3 → Output
  \foreach \i in {1,...,3} {
    \foreach \k in {1,2} {
      \draw[conn] (h3\i) -- (o\k);
    }
  }

  % Layer labels
  \node[layer label, above=0.4cm of i1]  {Input\\Layer};
  \node[layer label, above=0.4cm of h21] {Hidden Layers};
  \node[layer label, above=0.4cm of o1]  {Output\\Layer};

\end{tikzpicture}

    }
    \caption{Example of a feedforward neural network with three hidden layers.}
    \label{fig:FFNN_example}
\end{figure}

For example, taking in consideration the FFNN architecture shown in the figure \ref{fig:FFNN_example} with three layers, we can express the mapping as follows:
\begin{equation*}
    f(\vct{x}) = f^{(5)}(f^{(4)}(f^{(3)}(f^{(2)}(f^{(1)}(\vct{x})))))
\end{equation*}
\begin{itemize}
    \item The first layer $f^{(1)}(\vct{x})$ is usually referred to as the \textbf{input layer}, each node corresponds to a feature of the input data. Usually no
    computations are performed in this layer, as it simply serves as a conduit for the input data to be fed into the network;
    \item The layers $f^{(2)}$, $f^{(3)}$, and $f^{(4)}$ are called \textbf{hidden layers}, and they perform computations on the input data, applying a linear transformation followed by a non-linear activation function.
    The number of hidden layers and the number of neurons in each hidden layer depend on the architecture of the network and the complexity of the problem being solved.
    \item The last layer $f^{(5)}$ is called the \textbf{output layer}, and it produces the final output of the network. Depending on the type of problem,
    different activation functions can be used (e.g., linear for regression, softmax for classification)
    and different number of neurons can be used (e.g., one for binary classification and regression, one per class for multi-class classification).
\end{itemize}
The function composition can be described by a \textbf{directed acyclic graph} (DAG), where the nodes represent the functions (layers) and the edges represent the flow of data through the network. 
The term "feedforward" comes from the fact that the data flows in one direction, from the input layer to the output layer, without any cycles or loops in the graph.

\subsubsection{Difference from Linear Models}
Linear models are generally very easy to optimize (e.g., closed-form solutions or convex optimization), but they are limited in their expressiveness and cannot capture interactions between features or non-linear relationships in the data.
FFNNs allow to overcome the limitations of linear models by introducing non-linear activation functions and multiple layers of computation.

\paragraph{How can we improve Linear Models?}
FFNNs are not the only way to improve linear models, but they are one of the most powerful and widely used approaches.

In general, to boost the performance of linear models, we can apply a non-linear transformation to the input data 
$\vct{x} \rightarrow \phi(\vct{x})$, where $\phi$ is a non-linear function.
We can then apply the same linear model as before to the transformed data $\phi(\vct{x})$. 
The transformation $\phi$ can be implemented in various ways, such as:
\begin{itemize}
    \item \textbf{Kernel Methods}: map the input data into a high-dimensional feature space (possibly infinite-dimensional) where linear models can be applied.
    This has enough capacity to capture complex relationships in the data, but achieves poor generalization performance for highly variable target functions.
    \item \textbf{Feature Engineering}: manually design and select features that capture relevant information from the input data.
    This was the traditional approach before the rise of deep learning, in particular for computer vision tasks, but it is time-consuming and requires domain expertise.
    \item \textbf{Neural Networks}: learn the transformation $\phi$ from the data itself. 
    This shifts the engineering effort from designing features to designing the architecture of the network, but throwhs away the convexity and interpretability. 
\end{itemize}

\subsubsection{Training a FFNN}
To train a Neural Network, we need to optimize the parameters of the network (weights $\vct{\theta}$) to drive the output of the network $f(\vct{x}; \vct{\theta})$
to match the target variable $y$ for a given input $\vct{x}$.

The training process is nothing different from training any other machine learning model.
We need to define a loss function $\loss(f(\vct{x}; \vct{\theta}), y)$ that measures the discrepancy between the predicted output and the true target variable.

The largest difference with respect to linear models comes from the fact that most loss functions are non-convex, which removes the guarantee of finding a global minimum and makes the optimization process more challenging.

\subsubsection{Modelling Choices}
When designing a FFNN, we need to make several modelling choices, such as:
\begin{itemize}
    \item \textbf{Cost Function}: the choice of the loss function $\loss$ depends on the type of problem we are trying to solve (e.g., regression, classification) and the properties of the data.
    Each function has its own advantages and disadvantages, and the choice can significantly impact the performance of the model.
    \item \textbf{Form of Output}: the function applied in the output layer. Similarly to the cost function, the choice of the output units depends on the type of problem we are trying to solve;
    \item \textbf{Activation Function}: the non-linear function applied in the hidden layers. It's independent from the type of problem we are trying to solve, but it can significantly impact the
    performance of the model and the convergence during training.
    \item \textbf{Architecture}: the number of hidden layers, the number of neurons in each hidden layer, etc. 
    The architecture of the network can significantly impact the performance of the model and the convergence during training.
    \item \textbf{Optimizer}: how we update the parameters of the network during training.
    \item \textbf{Regularizer}: how we prevent the model from overfitting the training data and improve generalization performance. 
\end{itemize}